\begin{textsample}{NEG Dim 3 – pi – Score 1.00 – pi/pi\_em\_25.txt}  \label{ex:f3_neg_031}
SEÇÃO II FORMAÇÃO GERAL BÁSICA : Áreas de Conhecimento Formação Geral Básica, Alinhada à BNCC Por Formação Geral Básica ( FGB ), entende -se a parte do currículo que é comum a todos os estudantes e atende ao que orienta a BNCC ( BRASIL, 2018 ).
Neste sentido, contempla todas as áreas do conhecimento, e tem uma carga horária mais rígida ( podendo ser, no máximo, 1.800h para as três séries do Ensino Médio ), mas com a flexibilidade de cada Estado poder escolher a distribuição que melhor se ajuste a seu contexto.
Conforme modelo apresentado em tópico anterior, o Piauí optou por distribuir a carga horária da FGB de forma decrescente ( 800h / 1ª série ; 600h / 2ª série e 400h / 3ª série ).
Essa proposta se justifica pela possibilidade de adaptação dos estudantes ao novo, de forma gradual.
Inicia -se com o tempo maior para a parte comum, a fim de que o aprofundamento possa acontecer de forma mais intensa nas duas séries seguintes ( vide proposta dos itinerários formativos em capítulo posterior ).
Dada a distribuição de carga horária por ano, é preciso considerar também a organização das áreas do conhecimento de forma a proporcionar a integração entre os componentes – razão pela qual se optou por detalhar os componentes no quadro de habilidades de cada área – a fim de permitir que o espaço destinado a cada componente seja bem compreendido por todos os responsáveis pela implementação deste currículo.
Destaca -se que a Resolução nº 03/2018¹, em seu Art. 7º, ao propor que os currículos sejam pensados de forma a articular “ vivências e saberes dos estudantes e contribuindo para o desenvolvimento de suas identidades e condições cognitivas e socioemocionais ”, orienta que : > § 1º Atendidos todos os direitos e objetivos de aprendizagem instituídos na Base Nacional Comum Curricular ( BNCC ), as instituições e redes de ensino podem adotar formas de organização e propostas de progressão que julgarem pertinentes ao seu contexto, no exercício de sua autonomia, na construção de suas propostas curriculares e de suas identidades.
Sob essa perspectiva, o currículo que ora apresentamos tem como foco o protagonismo estudantil, compreendendo que esta é a “ espinha dorsal ” do Novo Ensino Médio², pois ao se tornarem protagonistas, os estudantes se tornam responsáveis por suas escolhas, e pela tomada de decisões sobre sua vida, a partir de um projeto de vida elaborado no início do ensino médio.
Com base nessa premissa, destacada em documento orientador para a implementação do Novo Ensino Médio³, é que as Redes e/ou os sistemas de ensino são orientados a repensarem seus currículos de forma a garantir as aprendizagens previstas para essa etapa da educação básica.
Nesse contexto, os currículos devem promover o desenvolvimento das competências gerais e específicas da BNCC.
Neste capítulo, apresenta -se o detalhamento do modelo adotado pela rede estadual para a parte comum do currículo, destacando -se, conforme Resolução 03/2018, que, na formação geral básica ( FGB ), competências e habilidades devem estar articuladas como um todo indissociável. > Art. 11.
A formação geral básica é composta por competências e habilidades previstas na Base Nacional Comum Curricular ( BNCC ) e articuladas como um todo indissociável, enriquecidas pelo contexto histórico, econômico, social, ambiental, cultural local, do mundo do trabalho e da prática social, e deverá ser organizada por áreas de conhecimento : > > I – linguagens e suas tecnologias ; > II – matemática e suas tecnologias ; > III – ciências da natureza e suas tecnologias ; > IV – ciências humanas e sociais aplicadas.
Contudo, para maior clareza e compreensão, ficam definidos alguns termos neste Documento Curricular de Referência ( DCR ) com base nas normativas legais que norteiam o currículo do ensino médio ( Lei nº 9.394/96 ; Lei nº 13.005/2014 ( PNE ) ; Lei nº 13.415/2017 ; Resolução MEC-CNE-CEB nº 3/2018 ( DCNEM ) ; Resolução MEC-CNE-CP nº 4/2018 ( BNCC-EM ) ; Resolução CEE/PI nº 124/2020 ; Resolução CNE/CP nº 1/2021 e Portaria MEC nº 1.432/2018 MEC ), assim como nos diversos documentos pedagógicos orientadores, mas com destaque para as definições do Art. 6º das DCNEM : - FORMAÇÃO INTEGRAL : é o desenvolvimento intencional dos aspectos físicos, cognitivos e socioemocionais do estudante por meio de processos educativos significativos que promovam a autonomia, o comportamento cidadão e o protagonismo na construção de seu projeto de vida ; - FORMAÇÃO GERAL BÁSICA : conjunto de competências e habilidades das áreas de conhecimento previstas na Base Nacional Comum Curricular ( BNCC ), que aprofundam e consolidam as aprendizagens essenciais do ensino fundamental, a compreensão de problemas complexos e a reflexão sobre soluções para eles ; - APRENDIZAGENS ESSENCIAIS : são definidas como conhecimentos, habilidades, atitudes, valores e a capacidade de os mobilizar, articular e integrar, expressando -se em competências.
As aprendizagens essenciais compõem o processo formativo de todos os educandos, como direito de pleno desenvolvimento da pessoa, seu preparo para o exercício da cidadania e qualificação para o trabalho. - UNIDADES CURRICULARES : elementos com carga horária pré-definida, formadas pelo conjunto de estratégias, cujo objetivo é desenvolver competências específicas, podendo ser organizadas em áreas de conhecimento, disciplinas, módulos, projetos, entre outras formas de oferta ; - ARRANJO CURRICULAR : seleção de competências que promovam o aprofundamento das aprendizagens essenciais demandadas pela natureza do respectivo itinerário formativo ; - DIVERSIFICAÇÃO : articulação dos saberes com o contexto histórico, econômico, social, ambiental, cultural local e do mundo do trabalho, contextualizando os conteúdos a cada situação, escola, município, estado, cultura, valores, articulando as dimensões do trabalho, da ciência, da tecnologia e da cultura ; - COMPETÊNCIAS : mobilização de conhecimentos, habilidades, atitudes e valores, para resolver demandas complexas da vida cotidiana, do pleno exercício da cidadania e do mundo do trabalho.
Para os efeitos da Resolução 3/2018 ( DCNEM ), com fundamento no caput do art. 35 -A e no § 1º do art. 36 da LDB, a expressão “ competências e habilidades ” deve ser considerada como equivalente à expressão “ direitos e objetivos de aprendizagem ” presente na Lei do Plano Nacional de Educação ( Lei nº 13.005/2014/PNE ) ; - HABILIDADES : conhecimentos em ação, com significado para a vida, expressas em práticas cognitivas, profissionais e socioemocionais, atitudes e valores continuamente mobilizados, articulados e integrados ; Objetos do conhecimento são os conteúdos, os assuntos abordados ao longo de cada componente ou unidade curricular, isto é, aquilo que será o meio para o desenvolvimento das competências e habilidades ; Objetivos de aprendizagem são as descrições concisas, claramente articuladas do que os estudantes devem saber e compreender, e do que sejam capazes de fazer numa etapa específica de sua escolaridade ( nesse caso, na etapa do Ensino Médio ).
São os objetivos em função dos resultados pretendidos para a aprendizagem que se pretende que os estudantes alcancem, de acordo com os objetos do conhecimento propostos e mediante o desenvolvimento das competências e habilidades.
Na área da educação, para melhor entendimento sobre competência, é importante dissecarmos a definição à luz da BNCC, segundo a qual é “ a mobilização de conhecimentos ( mobilizar conceitos, saberes, experiências, vivências ), habilidades ( práticas cognitivas e socioemocionais ), atitudes ( conduta, comportamento, posicionamento ) e valores ( princípios, preceitos, concepções, padrões ) para resolver demandas complexas da vida cotidiana, do pleno exercício da cidadania e do mundo do trabalho ”.
Portanto, trata -se de “ maneiras diferentes e intercambiáveis para designar algo comum, isto é, aquilo que todos os estudantes devem aprender na Educação Básica ( aprendizagens essenciais ), o que inclui tanto os saberes quanto a capacidade de mobilizá -los e aplicá -los ”.
Em outras palavras, refere -se à capacidade de mobilizar recursos, conhecimentos ou vivências/experiências para resolver questões/problemas do cotidiano, com desenvolvimento do pensamento científico, criativo, crítico e reflexivo, comunicação, autonomia, empatia, dentre outras habilidades, atitudes e valores imprescindíveis para o mundo contemporâneo.
Ao adotar esse enfoque, a BNCC indica que as decisões pedagógicas devem estar orientadas para o desenvolvimento de competências.
Por meio da indicação clara do que os alunos devem “ saber ” ( considerando a constituição de conhecimentos, habilidades, atitudes e valores ) e, sobretudo, do que devem “ saber fazer ” ( considerando a mobilização desses conhecimentos, habilidades, atitudes e valores para resolver demandas complexas da vida cotidiana, do pleno exercício da cidadania e do mundo do trabalho ), a explicitação das competências oferece referências para o fortalecimento de ações que assegurem as aprendizagens essenciais definidas na BNCC. ( BRASIL, 2018, p. 13 ) Assim, a formação integral postula o desenvolvimento de competências para aprender a aprender, saber lidar com a comunicação e com a informação cada vez mais disponível, atuar com discernimento, criticidade, tolerância e responsabilidade nos contextos das relações sociais, das culturas digitais, aplicar conhecimentos para resolver demandas do cotidiano, ter autonomia para tomar decisões, agir com proatividade, resolutividade, intermediar e solucionar conflitos, conviver e aprender com as diferenças e com as diversidades.
Com ênfase, Zabala e Arnau ( 2014 ) aduzem que a competência deve detectar o que cada pessoa precisa para resolver questões cotidianas ao longo da vida : > A competência, no âmbito da educação escolar, deve identificar o que qualquer pessoa necessita para responder aos problemas aos quais será exposta ao longo da vida.
Portanto, a competência consistirá na intervenção eficaz nos diferentes âmbitos da vida, mediante ações nas quais se mobilizam, ao mesmo tempo e de maneira inter-relacionada, componentes atitudinais, procedimentais e conceituais ( ZABALA e ARNAU, 2014, p. 12 ).
Neste sentido, um currículo baseado no desenvolvimento de competências é aquele que reflete a formação em aprendizagens traduzidas na capacidade de serem aplicadas em contextos reais, cotidianos.
Portanto, a essência das competências é a sua funcionalidade nas situações reais, ou seja, sua efetiva aplicabilidade ( mobilização de conhecimentos ) diante de qualquer situação nova ou conhecida.
Quanto às habilidades, estas são os “ conhecimentos em ação, expressas em práticas profissionais e socioemocionais, atitudes e valores continuamente mobilizados, articulados e integrados ”.
São os conhecimentos em movimento ( saber fazer ) necessários para o pleno desenvolvimento das competências.
Todavia, o desenvolvimento de competências passa pela articulação de várias habilidades que são organizadas na BNCC de maneira progressiva, ou seja, das mais simples para as mais complexas.
De acordo com a Resolução nº 3/2018 ( DCNEM ) em seu Art. 17, § 7º, > As áreas do conhecimento podem ser organizadas em unidades curriculares, competências e habilidades, unidades de estudo, módulos, atividades, práticas e projetos contextualizados ou diversamente articuladores de saberes, desenvolvimento transversal ou transdisciplinar de temas ou outras formas de organização.
Seguindo essa orientação, a estrutura curricular desenhada na Formação Geral Básica para cada área do conhecimento, as competências, as habilidades, os objetos do conhecimento e os objetivos de aprendizagem estão organizados considerando conceitos estruturantes, temas integrados, situações problemas, situações de aprendizagem, unidades temáticas, dentre outras estratégias pedagógicas.
Tudo numa escala de complexidade e integralidade, observando a progressão da aprendizagem, partindo dos conhecimentos gerais para os específicos.
Nessa organização, os objetos do conhecimento são as expressões dos componentes curriculares ( e dos estudos e práticas ) que, em conjunto, potencializam o desenvolvimento das habilidades e competências das áreas do conhecimento.
No entanto, é importante enfatizar, que cada escola, com base ( e a partir ) do seu Projeto Político Pedagógico deverá adotar as estratégias pedagógicas que mais se adequem à sua realidade de acordo com sua autonomia e realidade, bem como consoante às diretrizes nacionais e estaduais, como também dos regramentos da rede de ensino à qual é vinculada.
É importante enfatizar que, enquanto as habilidades da Formação Geral Básica são definidas a partir das Competências Gerais e dos objetos do conhecimento, nos Itinerários Formativos, os objetos do conhecimento são definidos em função das habilidades gerais e específicas a serem desenvolvidas no âmbito dos quatro eixos estruturantes.
A forma como se organizam as competências, as habilidades, os objetos de conhecimento e os objetivos de aprendizagem por cada área do conhecimento com seus respectivos componentes curriculares, será detalhada em tópicos seguintes deste DCR.
Ademais, neste mesmo documento, contém orientação para que não haja prejuízo da integração e articulação entre as áreas do conhecimento.
Por essa razão, optamos por uma distribuição de carga horária que fosse equivalente para cada área, de modo a não haver sobreposição de um componente sobre os outros, nem de uma área sobre as outras.
A organização por áreas do conhecimento é uma das possibilidades de articular o currículo de forma interdisciplinar.
As áreas do conhecimento favorecem a comunicação entre os conhecimentos e saberes dos diferentes componentes curriculares, mas permitem que os procedimentos e conceitos próprios de cada componente sejam preservados.
Contudo, é importante lembrar que, para que a integração curricular ocorra de forma fluente para o aprendizado do aluno, é imprescindível que o planejamento seja realizado de forma integrada e que haja contextualização, interdisciplinaridade e transdisciplinaridade.
Outro aspecto que merece atenção neste texto é a presença dos componentes em cada área do conhecimento.
A Resolução diz que deve haver estudos e práticas de : I – língua portuguesa ; II – matemática ; III – conhecimento do mundo físico e natural [ … ] ; IV – arte ; V – educação física ; VI – história do Brasil e do mundo ; VII – história e cultura afro-brasileira e indígena ; VIII – sociologia e filosofia ; IX – língua inglesa.
No entanto, este currículo denomina e detalha cada componente curricular em quadro de competências e habilidades, inclusive, com orientações de objeto do conhecimento a ser trabalhado.
Essa escolha se justifica pelo entendimento de que as escolas da Rede ainda não têm maturidade para vivenciar uma mudança tão brusca ( considerando que hoje se tem um currículo conteudista e por componente curricular ) e que, mesmo com todo o processo formativo pelo qual passarão as equipes das escolas, ainda é preciso cautela na implementação do “ novo ”.
Sob essa ótica, evidencia -se que, além das orientações dos normativos nacionais, a exemplo da Resolução 03/2018, aqui citada, para a construção do currículo, foram consideradas as potencialidades do território piauiense, conforme descritas nas propostas das quatro áreas de conhecimento.
Em conformidade com os fundamentos pedagógicos da BNCC, este DCR está estruturado de modo a explicitar as competências que devem ser desenvolvidas ao longo de todo o Ensino Médio, em cada área do conhecimento e seus respectivos componentes curriculares, como expressão dos direitos de aprendizagem e desenvolvimento de todos os estudantes piauienses.
Por fim, salienta -se que, nos tópicos seguintes, encontram -se as indicações para o atendimento da garantia do caráter interdisciplinar e transdisciplinar entre as áreas, bem como orientações para efetivação da proposta apresentada neste documento curricular.
Antes, porém, discorremos sobre a progressão das aprendizagens de uma etapa à outra, mais especificamente do Ensino Fundamental para o Médio. -- - ¹ Disponível em https://www.in.gov.br/materia/-/asset\_publisher/Kujrw0TZC2Mb/content/id/51281622.
Acesso em 22/10/2020. ² Nos manuais de orientação sobre o Novo Ensino Médio, a espinha dorsal é o protagonismo juvenil.
Entretanto, por razões já elencadas em documento introdutório deste currículo, utilizamos a expressão protagonismo estudantil, com vistas a atender a todos os sujeitos da referida etapa. ³ Cf.
Guia de Implementação do Novo Ensino Médio ( MEC/Consed, 2019 ).
Disponível em http://novoensinomedio.mec.gov.br/resources/downloads/pdf/Guia.pdf.

% matched lemmas: 
\end{textsample}
